\ssscp{Innovations among the numerals}{07-04-02-01}
\trf{07-14} and \trf{07-15} show developments in the lexicon
in the shape of numerals 1--9
and the etymon for ‘ten’ that distinguish \tsc{Sula-Buru} from \tsc{Ambon-Seram}.
The unity across all \tsc{Sula-Buru}
and \tsc{Ambon-Seram} languages of **utum ‘hundred’ is shown in \trf{07-06}.
Most \tsc{Seram-Tanimbar-Bomberai} languages have reflexes of
PMP *Ratus ‘hundred’ and not **utum, and most also share **butu ‘ten’ with \tsc{Ambon-Seram} languages.
In the larger context most \tsc{Timor-Babar} languages of the
\tsc{Southwest Maluku} node also have ‘ten’ from **butu (\trf{04-46} \prf{tab:04-46}).
In the \tsc{Bima-Lembata} subgroup,
there has been a shift in reflexes of *butu in \tsc{Central Flores}
and \tsc{Flores-Lembata} languages to mean ‘four’
(\citealp[367]{Fricke2019}, \trf{03-36} \prf{tab:03-36}).

\begin{table}[p]\centering\tcp{Central Maluku numerals 1--5}{07-14}
\begin{adjustbox}{totalheight=\textheight}
	%\st{2pt}
	\begin{tabular}{llllll}\lsptoprule
*gloss	&	one	&	two	&	three	&	four	&	five\\
PMP	&	*əsa, *isa, *asa	&	*duha	&	*təlu	&	*əpat	&	*lima\\ \midrule
%\tsc{Sula-Buru}\\ \midrule
Mangole	&	{\hr}gia, ga-ia	&	{\hr}guu, ga-uu	&	{\hr}ga-telu	&	\cg{(ga-dia)}	&	{\hr}ga-limo\\
Sanana	&	{\hr}hia	&	{\hr}ga-hu	&	{\hr}ga-tel	&	\cg{(ga-reha)}	&	{\hr}ga-lima\\
Lisela	&	{\hr}m-sia-n, saa	&	{\hr}rua	&	{\hr}telo	&	{\hr}paa	&	{\hr}lima\\
Buru	&	{\hr}em-sia-n, saa	&	{\hr}rua	&	{\hr}telo	&	{\hr}paa	&	{\hr}lima\\
Ambelau	&	{\hr}esbii	&	{\hr}lua	&	{\hr}relo	&	{\hr}faa	&	{\hr}lima\\ \midrule
%\tsc{Ambon-Seram}\\ \midrule
Kayeli	&	{\hr}sileʔ {\tl} sile-i	&	{\hr}luaʔ	&	{\hr}telo	&	{\hr}haa (<SB)	&	{\hr}lima\\
Boano	&		&	{\hr}lua	&	{\hr}tenu	&		&	\\
Piru	&	{\hr}saa, wasa	&	{\hr}lua	&	{\hr}telu	&	{\hr}ata	&	{\hr}lima\\
Larike	&	{\hr}-sa, sane	&	{\hr}dua	&	{\hr}tidu	&	{\hr}ati	&	{\hr}dima\\
Asilulu	&	{\hr}sa	&	{\hr}lua	&	{\hr}telu	&	{\hr}ata	&	{\hr}lima\\
Paulohi	&	{\hr}sae	&	{\hr}ʀua	&	{\hr}toʀu	&	\cg{(fale)}	&	{\hr}ʀima\\
Amahai	&	\cg{(kane)}	&	{\hr}ruwa	&	{\hr}oru	&	{\hr}haʔa	&	{\hr}rima\\
Kaibobo	&	{\hr}isa	&	{\hr}rua	&	{\hr}toru	&	{\hr}ata	&	{\hr}lima\\
Kamarian	&	{\hr}isa	&	{\hr}rua	&	{\hr}teru	&	{\hr}haʔa	&	{\hr}rima\\
Haruku (Karu)	&	{\hr}isa	&	{\hr}rua	&	{\hr}toru	&	{\hr}haʔa	&	{\hr}rima\\
Haruku	&	{\hr}sa, sane	&	{\hr}rua	&	{\hr}toru	&	{\hr}haʔa	&	{\hr}rima\\
Saparua (Latu)	&		&		&	{\hr}todu	&	{\hr}haʔa	&	\\
Saparua (Hatawano)	&	{\hr}isa	&	{\hr}lua	&	{\hr}tolu	&	{\hr}haʔa	&	{\hr}lima\\
Nusalaut	&	{\hr}isa	&	{\hr}rua	&	{\hr}toru, ʔoru	&	{\hr}haʔa, hata	&	{\hr}rima\\
Sepa	&	{\hr}isa, san	&	{\hr}rua	&	{\hr}toru	&	\cg{(hale)}	&	{\hr}rima\\
Tamilou	&	{\hr}isa	&	{\hr}lua	&	{\hr}tou	&	\cg{(hae)}	&	{\hr}lima\\
Tehoru	&	{\hr}san	&	{\hr}lua	&	{\hr}toi	&	\cg{(fale)}	&	{\hr}lima\\
N.Wemale	&	{\hr}sai	&	{\hr}lua	&	{\hr}telu	&	\cg{(hale)}	&	{\hr}lima\\
N. Alune	&	{\hr}esa	&	{\hr}lua	&	{\hr}telu	&	{\hr}ata	&	{\hr}lima\\
W. Alune	&	{\hr}esa	&	{\hr}dua	&	{\hr}tedu	&	{\hr}ata	&	{\hr}dima\\
Yalahatan	&	{\hr}isa, sa	&	{\hr}rua	&	{\hr}telu	&	{\hr}ata	&	{\hr}lima\\
S. Nuaulu	&	{\hr}isa, osa	&	{\hr}ua	&	{\hr}tonu	&	{\hr}ate	&	{\hr}nima\\
Manusela	&	{\hr}esa	&	{\hr}hua	&	{\hr}tolu	&	{\hr}hate	&	{\hr}lima\\
Liana-Seti	&	{\hr}esa	&	{\hr}en-lua	&	{\hr}en-tol	&	{\hr}en-hata	&	{\hr}en-lima\\ \midrule
Banda	&	{\hr}sa	&	{\hr}ruo	&	{\hr}telu, tel	&	{\hr}at	&	{\hr}limo\\ \midrule
%\tsc{Seram-Tanimbar-Bomberai}\\ \midrule
Bobot	&	{\hr}san	&	{\hr}lua	&	{\hr}tolw	&	{\hr}fet	&	{\hr}lima\\
Masiwang	&	{\hr}sa, ya	&	{\hr}lua	&	{\hr}toli	&	{\hr}faat	&	{\hr}lima\\
Geser	&	{\hr}sa	&	{\hr}roti	&	{\hr}tolu	&	{\hr}fat	&	{\hr}lim\\
Gorom	&	{\hr}sa	&	{\hr}roti	&	{\hr}tolu	&	{\hr}hat	&	{\hr}lim, lium\\
Watubela	&	{\hr}ha	&	{\hr}lua	&	{\hr}tolu	&	{\hr}fata	&	{\hr}lima\\
Kowiai	&	{\hr}sa, \cg{(isamos)}	&	{\hr}ruwet	&	{\hr}tor	&	{\hr}ɸat	&	{\hr}rim\\
Teor	&		&	‹rúa›	&	‹teili›	&	‹faht›	&	‹lima›\\
Kur	&	{\hr}ha	&	{\hr}rua	&	{\hr}tel	&	‹paät›	&	‹lima›\\
Yamdena	&	{\hr}sa, \cg{(lese)}	&	{\hr}du, due	&	{\hr}tel, tely	&	{\hr}fat	&	{\hr}lim\\
Fordata	&	{\hr}saa	&	{\hr}rua	&	{\hr}telu	&	{\hr}faʔat	&	{\hr}lima\\
Uruangnirin	&	{\hr}sa	&	{\hr}nua	&	{\hr}teni	&	{\hr}fat	&	{\hr}nima\\
Onin	&	{\hr}sa	&	{\hr}nua	&	{\hr}teni	&	{\hr}fat	&	{\hr}nima\\
Sekar	&	{\hr}isa	&	{\hr}nua	&	{\hr}teni	&	{\hr}fat	&	{\hr}nima\\
Arguni	&	{\hr}syá	&	{\hr}ru	&	{\hr}táòr	&	{\hr}fat	&	{\hr}rím\\
Bedoanas	&	{\hr}sa	&	{\hr}ruyu	&	{\hr}taur	&	{\hr}fat	&	{\hr}rim\\
Erokwanas	&	{\hr}sa	&	{\hr}ru	&	{\hr}tour	&	{\hr}fat	&	{\hr}rim\\
Irarutu	&	{\hr}esu	&	{\hr}ru, rifu	&	{\hr}tur	&	\cg{(gigiti)}	&	\cg{(fradfid)}\\
		\lspbottomrule
	\end{tabular}
\end{adjustbox}
\end{table}

\begin{table}[p]\centering\tcp{Central Maluku numerals 6--10 }{07-15}
\begin{adjustbox}{totalheight=\textheight}
	\begin{threeparttable} \st{2.5pt}
	\begin{tabular}{llllll@{ }l}\lsptoprule
*gloss	&	six	&	seven	&	eight	&	nine	&	ten &ten\\
PMP	&	*ənəm	&	*pitu	&	*walu	&	*siwa	&	*puluq, &**butu\\ \midrule
%\tsc{Sula-Buru}\\ \midrule
Mangole	&	{\hr}ga-nee	&	{\hr}ga-pitu	&	\cg{(ga-ta-hua)}	&	\cg{(ga-ta-sia)}\su{†}	&	{\hr}poo &\\
Sanana	&	{\hr}ga-nei	&	{\hr}ga-pit	&	\cg{(ga-ta-hua)}	&	\cg{(ga-ta-sia)}\su{†}	&	{\hr}poa {\tl} poo&\\
Lisela	&	{\hr}nee	&	{\hr}pito	&	\cg{(t-rua)}	&	\cg{(ʧia /t-sia/)}\su{†}	&	{\hr}polo {\tl} pulu&\\
Buru	&	{\hr}nee	&	{\hr}pito	&	\cg{(et-rua)}	&	\cg{(eʧia /et-sia/}\su{†}	&	{\hr}polo&\\
Ambelau	&	{\hr}nee	&	{\hr}fiʧo {\tl} fito	&	{\hr}βalo	&	{\hr}siβa	&	&{\hr}poro (<**mbutu)\\ \midrule
%\tsc{Ambon-Seram}\\ \midrule
Kayeli	&	{\hr}nee (<SB)	&	{\hr}hitoʔ{\tl}itoʔ	&	{\hr}waloʔ	&	{\hr}siwaʔ	&	&{\hr}botoʔ\\
Piru	&	{\hr}nena	&	{\hr}itu	&	{\hr}walu	&	{\hr}siwa	&&	{\hr}hutu\\
Larike	&	{\hr}nene	&	{\hr}itu	&	{\hr}wadu	&	{\hr}siwa	&&	{\hr}husa (<**hut-sa)\\
Asilulu	&	{\hr}nena	&	{\hr}itu	&	{\hr}walu	&	{\hr}siwa	&&	{\hr}hutu (besi)\\
Paulohi	&	{\hr}noyo	&	{\hr}fitu	&	{\hr}waʀu	&	{\hr}siwa	&&	{\hr}hutu sae\\
Amahai	&	{\hr}noʔo	&	{\hr}hitu	&	{\hr}waru	&	{\hr}siwa	&&	{\hr}huʔu tae\\
Kaibobo	&	{\hr}noʔo	&	{\hr}itu	&	{\hr}aru	&	{\hr}sia	&&	{\hr}hutu\\
Kamarian	&	{\hr}neʔe	&	{\hr}itu	&	{\hr}waru	&	{\hr}siwa	&&	{\hr}hutu\\
Haruku	&	{\hr}noʔo	&	{\hr}itu	&	{\hr}waro	&	{\hr}siwa	&&	{\hr}huʔu sa\\
Haruku (Karu)	&	{\hr}noʔo	&	{\hr}itu	&	{\hr}waru	&	{\hr}siwa	&&	{\hr}hutu\\
Saparua (Hatawano)	&	{\hr}noʔo	&	{\hr}hitu	&	{\hr}walu	&	{\hr}siwa	&&	{\hr}hutu\\
Nusalaut	&	{\hr}noʔo	&	{\hr}hitu, hiu	&	{\hr}walu	&	{\hr}siwa	&&	{\hr}huʔu sai\\
Sepa	&	{\hr}noo	&	{\hr}hitu	&	{\hr}waru	&	{\hr}siwa	&&	{\hr}hutu sai\\
Tamilou	&	{\hr}noe	&	{\hr}hitu	&	{\hr}waʤo	&	{\hr}siwa	&&	{\hr}hutu sai\\
Tehoru	&	{\hr}noe	&	{\hr}fitu	&	{\hr}waʤu	&	{\hr}siwa	&&	{\hr}hutu san\\
N. Wemale	&	{\hr}tiniane	&	{\hr}hitu	&	{\hr}walu	&	{\hr}siwa	&&	{\hr}putu sa\\
N. Alune	&	{\hr}nee	&	{\hr}itu	&	{\hr}kwalu	&	{\hr}sikwa	&&	{\hr}butu sa\\
W. Alune	&	{\hr}neʔe	&	{\hr}itu	&	{\hr}kwadu	&	{\hr}sikwa	&&	{\hr}butu\\
Yalahatan	&	{\hr}nome	&	{\hr}itu	&	{\hr}walu	&	{\hr}siwa	&&	{\hr}hutu-sa\\
S. Nuaulu	&	{\hr}nome	&	{\hr}itu	&	{\hr}wanu	&		&&	{\hr}hutu-sa\\
Manusela	&	{\hr}nome	&	{\hr}hitu	&	{\hr}walu	&	{\hr}siwa	&&	{\hr}hutu sa\\
Liana-Seti	&	{\hr}en-noi	&	{\hr}en-hit	&	{\hr}wal	&		&&	{\hr}fut-sa\\ \midrule
Banda	&	{\hr}namu	&	{\hr}itu	&	{\hr}walu	&	{\hr}siwa	&&	{\hr}futusa\\ \midrule
%\tsc{Seram-Tanimbar-Bomberai}\\ \midrule
Bobot	&		&		&	{\hr}walw	&		&	\\
Masiwang	&	{\hr}noon	&	{\hr}fiti	&	{\hr}wali	&	{\hr}sifa	&&	{\hr}fuʧa\\
Geser	&	{\hr}onan	&	{\hr}fitu	&	{\hr}alu	&	{\hr}sia	&&	{\hr}utsa [uʧa]\\
Gorom	&	{\hr}onan, onaŋ	&	{\hr}hitu	&	{\hr}alu	&	{\hr}sia	&&	{\hr}utsa [uʧa]\\
Watubela	&	{\hr}onon	&		&		&		&&	\\
Kowiai	&	\cg{(rim isamos)}	&	\cg{(rim wet)}	&	\cg{(rim tor)}	&	\cg{(rim bat)}	&&	{\hr}ɸutsa [ɸuʧa]\\
Teor	&	‹nem›	&	‹fit›	&	‹wal›	&	‹siwer›	&&	‹hutá›\\
Kur	&	‹nèèn›	&	‹piet›	&	‹wah›	&	‹siwah›	&	\mc{2}{l}{‹sapoel› < Mly?}\\
Yamdena	&	{\hr}nem	&	{\hr}itw	&	{\hr}walu	&	{\hr}siw	&&	{\hr}buty\\
Fordata	&	{\hr}nean	&	{\hr}fitu	&		&		&&	{\hr}vutu\\
Uruangnirin	&	{\hr}nem	&	\cg{(taransa)}	&	\cg{(terinua)}	&	\cg{(saputi)}	&&	{\hr}puʧa\\
Onin	&	{\hr}nem	&	\cg{(tara sa)}	&	\cg{(tara nua)}	&	\cg{(saputi)}	&&	{\hr}pusua\\
Sekar	&	{\hr}nem	&	\cg{(buteras)}	&	\cg{(beterua)}	&	\cg{(masbuti)}	&&	{\hr}buswa\\
Arguni	&	{\hr}ánem	&	\cg{(nambatu)}	&	\cg{(nafu)}	&	{\hr}neswá	&	{\hr}sambure&\\
Bedoanas	&	{\hr}unam	&	\cg{(nambatu)}	&	\cg{(nafu)}	&	{\hr}naswa	&	{\hr}sambura&\\
Erokwanas	&	{\hr}aɲam	&	\cg{(nambatu)}	&	\cg{(nafu)}	&	{\hr}naswa	&	{\hr}sambura&\\
Irarutu	&	\cg{(tresu)}	&	\cg{(treru)}	&	\cg{(tretur)}	&	\cg{(tregigiti)}	&	\cg{(fradru)}&\\
		\lspbottomrule
	\end{tabular}
		\begin{tablenotes}\tnsize
			\item[†] {Despite the superficial similarity,
								the \tsc{Nuclear Sula-Buru} words are not from *siwa but are
								phrases meaning ‘minus-one’ formed with reflexes of *isa > **sia ‘one’.
								Compare the words for ‘eight’ which are ‘minus-two’ (see discussion
								in \srf{08-04}).}
		\end{tablenotes}
	\end{threeparttable}
\end{adjustbox}
\end{table}
